\section{Corrections and windows}\label{sec:corrections}

This section treats arbitrary sets and finite windows of targets. The undershoot of an arbitrary set at any level is that of $I_h$ plus an explicit correction, and the undershoot over a window has a closed form in the corners of Section~\ref{sec:noncm}. With $D_E$ and $F_h$ as in Section~\ref{sec:noncm} put
\[
o_2(y)=\min\{j\ge0:y+j\in S_2\}.
\]
An \emph{equality level} is an $h$ with $hA_0=I_h$: every $h$ for Cohen--Macaulay sets, every $h\ge\tau$ otherwise.

\subsection{Undershoot as overshoot}\label{subsec:overshoot-form}

\begin{proposition}[Head and reflected tail]\label{prop:overshoot-form}
Let $A_0$ be arbitrary and $h\ge0$.
\begin{enumerate}
\item[(a)] $D_{I_h}(n)=D_{S_1}(n)$ for $0\le n\le hd-c_2$.
\item[(b)] If $x\in I_h$ and $x\ge c_1$, then $D_{I_h}(n)=o_2(hd-n)$ for $x\le n\le hd$. One may take $x=c_1$ when $hd-c_1\in S_2$, in particular when $c_1=0$ or $hd\ge c_1+c_2$.
\item[(c)] If $hd\ge c_1>0$ and $N_h=\max(I_h\cap[0,c_1-1])$, then for $c_1\le n\le hd$, $D_{I_h}(n)=o_2(hd-n)$ if $o_2(hd-n)\le n-c_1$ and $D_{I_h}(n)=n-N_h$ otherwise.
\item[(d)] If $hd\ge c_1+c_2-1$, (a) together with (b) or (c) determines $D_{I_h}$ on $[0,hd]$.
\end{enumerate}
At an equality level $D_h=D_{I_h}$.
\end{proposition}

\begin{proof}
(a) For $n'\le n\le hd-c_2$, $hd-n'\ge c_2$, so $n'\in I_h$ if and only if $n'\in S_1$. (b), (c) For $c_1\le n'\le n$, $n'\in S_1$, so $n'\in I_h$ if and only if $hd-n'\in S_2$. With $j=o_2(hd-n)$: if $j\le n-c_1$, then $n-j$ is the largest element of $I_h$ in $[c_1,n]$; otherwise there is none and the predecessor is $N_h$ (and $0\in I_h$). If $x\in I_h$, $x\ge c_1$ and $n\ge x$, then $j\le n-x$, since $(hd-n)+(n-x)=hd-x\in S_2$. If $c_1=0$, $hd-c_1=hd\in S_2$. (d) When $hd\ge c_1+c_2-1$ we have $hd-c_2\ge c_1-1$, so the range $[0,hd-c_2]$ of (a) and the range $[c_1,hd]$ of (b) and (c) together cover $[0,hd]$.
\end{proof}

So above the anchor the undershoot is the overshoot of the reflected target into $S_2$; summed over $[0,hd]$ this is Proposition~\ref{prop:overshoot}. For an arbitrary set the actual undershoot differs from that of $I_h$ by an explicit correction.

\begin{proposition}[Correction]\label{prop:correction}
For $z\in F_h$ let $z^+=\min(I_h\cap(z,hd])$ and $z^-=\max(hA_0\cap[0,z))$. For every $W\subseteq[0,hd]$,
\[
\sum_{n\in W}D_h(n)=\sum_{n\in W}D_{I_h}(n)+\Psi_h(W),\qquad \Psi_h(W)=\sum_{z\in F_h}\#\bigl(W\cap[z,z^+)\bigr)(z-z^-).
\]
So $\Psi_h\ge0$, and $\Psi_h([0,hd])>0$ if and only if $h$ is a failing level.
\end{proposition}

\begin{proof}
Both $z^\pm$ exist as $0,hd\in hA_0$. For a target $n$ let $z=\max(I_h\cap[0,n])$. If $z\notin F_h$ it is also the predecessor in $hA_0\subseteq I_h$. If $z\in F_h$, then $n\in[z,z^+)$, $hA_0$ has no element in $[z,n]$, and the predecessor is $z^-$, adding $z-z^-$. The cells $[z,z^+)$ are disjoint.
\end{proof}

For weighted targets replace the cardinality by the weight of $W\cap[z,z^+)$.

\subsection{Window formulas}\label{subsec:overlap}

For arbitrary sets the correction on a window has a finite closed form in the corners of Proposition~\ref{prop:failing-interval}. Let $C_s$ be the set of minimal points of $\mathcal C_s$, so that $kd-s\in hA_0$ if and only if $\kappa\le k\le h-\mu$ for some $(\kappa,\mu)\in C_s$. These corners are, in other coordinates, the strict records of the least number of parts in the class of $-s$ that Chen and Yang use, and for a single position Lemma~\ref{lem:hinge} is the positive-part count behind their exact formula for $|hA_0|$ \cite[Lemma~2, Thm.~1]{ChenYang2026}; the window clipping, the union over consecutive positions and the passage to the undershoot and its correction are what Theorem~\ref{thm:hinge} adds. For $1\le s\le t<d$ let $C^B_{s,t}$ be the minimal elements of $\bigcup_{r=s}^tC_r$ and $C^I_{s,t}$ those of $\{(u_r,v_r):s\le r\le t\}$, listed with $\kappa_1<\dots<\kappa_J$ and $\mu_1>\dots>\mu_J$. Each corner $(\kappa_j,\mu_j)$ makes a position present in the blocks $k\in[\kappa_j,h-\mu_j]$, and the window keeps the blocks in $[\ell_0,q_0]$. The function $H_C$ below counts the blocks of $[\ell_0,q_0]$ covered by at least one of these intervals: it adds their clipped lengths and subtracts the overlap of each interval with the one before it, which is all of its overlap with the earlier ones, since both endpoint sequences increase. With $[x]_+=\max(x,0)$ put
\begin{multline*}
H_C(h;\ell_0,q_0)=\sum_{j=1}^{J}\bigl[\min(q_0,h-\mu_j)-\max(\ell_0,\kappa_j)+1\bigr]_+\\
-\sum_{j=2}^{J}\bigl[\min(q_0,h-\mu_{j-1})-\max(\ell_0,\kappa_j)+1\bigr]_+ .
\end{multline*}

\begin{lemma}\label{lem:hinge}
$H_C(h;\ell_0,q_0)$ is the number of integers in $[\ell_0,q_0]\cap\bigcup_j[\kappa_j,h-\mu_j]$.
\end{lemma}

\begin{proof}
The intervals $I_j=[\max(\ell_0,\kappa_j),\min(q_0,h-\mu_j)]$ have nondecreasing endpoints, so a point of $I_j\cap I_i$, $i<j-1$, lies in $I_{j-1}$, and $I_j\cap(I_1\cup\dots\cup I_{j-1})=I_j\cap I_{j-1}$, of size given by the $j$-th subtracted term (also for empty intervals).
\end{proof}

\begin{theorem}[Hinge formulas]\label{thm:hinge}
Let $T_-\le T_+$ be integers in $[0,hd]$, and put $\ell_s=\max(1,\lceil(T_-+s)/d\rceil)$ and $q_s=\min(h,\lfloor(T_++s)/d\rfloor)$. Then
\begin{gather*}
\sum_{T=T_-}^{T_+}D_h(T)=\sum_{1\le s\le t<d}\Bigl([q_s-\ell_s+1]_+-H_{C^B_{s,t}}(h;\ell_s,q_s)\Bigr),\\
\Psi_h([T_-,T_+])=\sum_{1\le s\le t<d}\bigl(H_{C^I_{s,t}}(h;\ell_s,q_s)-H_{C^B_{s,t}}(h;\ell_s,q_s)\bigr).
\end{gather*}
For $0<c<d$, $c\notin A_0$, with corners $C'$ for $A_0\cup\{c\}$, the total undershoot of $A_0\cup\{c\}$ at level $h$ on $[T_-,T_+]$ is smaller than that of $A_0$ at level $h+1$ if and only if $\sum_{s\le t}H_{C'^B_{s,t}}(h;\ell_s,q_s)>\sum_{s\le t}H_{C^B_{s,t}}(h+1;\ell_s,q_s)$, and they tie exactly when the sums are equal.
\end{theorem}

\begin{proof}
Multiples $kd$, $k\le h$, lie in $hA_0$, so as in Proposition~\ref{prop:counting}(a), $D_h(kd-s)$ counts the $t\ge s$ such that no position of $[s,t]$ is present in block $k$. The target $kd-s$ lies in $[T_-,T_+]$ exactly for $\ell_s\le k\le q_s$. Position $r$ is present in block $k$ exactly for $k\in[\kappa,h-\mu]$, $(\kappa,\mu)\in C_r$, and dominated corners give smaller intervals; Lemma~\ref{lem:hinge} counts the blocks in which some position of $[s,t]$ is present. For $I_h$ the corners are the $(u_r,v_r)$ (Lemma~\ref{lem:residue}). For the comparison, $T_+\le hd$ and $s<d$ give $\lfloor(T_++s)/d\rfloor\le h$, so $\ell_s,q_s$ are the same at levels $h$ and $h+1$ and the first terms cancel.
\end{proof}

The corners depend neither on $h$ nor on the window, so finitely many integers per residue class give the undershoot over every window at every level.
For $[T_-,T_+]=[0,hd]$ every term is $[h+1-\kappa-\mu]_+$ with fixed corners, so $\Psi_h([0,hd])$ is piecewise affine in $h$ and vanishes for $h\ge\tau$.
